A análise dos dados apresenta algumas ameaças à validade que devem ser consideradas na interpretação dos resultados. Inicialmente, métricas baseadas na presença do aluno, como o tempo médio real de sessão, podem superestimar o engajamento devido à possibilidade de ``presença fantasma'', quando o estudante mantém a aba aberta sem participação cognitiva efetiva. Além disso, a dependência da frequência de emissão de \textit{traces} pela extensão representa uma limitação: intervalos longos entre registros podem tanto refletir problemas técnicos quanto desengajamento não capturado, afetando múltiplas métricas.

A taxa de presença também sofre ameaças específicas, como o truncamento em 100\%, que pode mascarar engajamentos acima do previsto, bem como distorções ao comparar o tempo médio individual com a duração total da aula, especialmente em turmas com grande variabilidade de participação. Em relação às métricas de distração, a classificação binária aplicada pela função de detecção pode não captar nuances de navegação educacional, e o método de atribuição do estado de distração a todo o intervalo entre dois \textit{traces} tende a superestimar tempos caso a mudança de estado ocorra dentro do período amostral. A sensibilidade dessa métrica à frequência de coleta também pode levar tanto à perda de transições rápidas quanto à captura de flutuações irrelevantes.

As medidas de frequência de distração e de troca de aba compartilham desafios relacionados à granularidade temporal e ao contexto de navegação: transições rápidas podem passar despercebidas, enquanto mudanças de URL dentro de um mesmo site podem inflar artificialmente o indicador de multitarefa. Além disso, essas métricas não distinguem trocas relacionadas à atividade da aula, como consulta a materiais complementares, de distrações reais. A fragmentação do foco também apresenta limitações metodológicas, como a contagem de retornos breves à aba como novos segmentos e a ausência de um limiar mínimo de duração para que um segmento seja considerado válido.

As métricas de engajamento audiovisual --- uso de câmera e microfone --- refletem ameaças adicionais. A atribuição temporal baseada no estado do \textit{trace} anterior pode superestimar o tempo de ativação de dispositivos, e a mera ativação desses recursos não garante engajamento cognitivo. Fatores externos, como normas institucionais, privacidade doméstica e limitações tecnológicas, também podem interferir no comportamento dos alunos sem refletir seu real envolvimento.

As métricas de tendência temporal dependem de suposições específicas, como a linearidade da distribuição das atividades durante a aula e a adequação de limiares fixos (por exemplo, a queda abaixo de 50\% do pico). Com baixa frequência de \textit{traces}, registros próximos às fronteiras dos quartis podem ser atribuídos de forma equivocada. Picos e quedas também podem ser afetados por momentos estruturais da aula (como pausas), que não refletem engajamento pedagógico real.

O \textit{score} geral de engajamento tem como ameaça a arbitrariedade dos pesos atribuídos a cada componente e ao risco de reduzir comportamentos complexos a uma métrica unidimensional.

Por fim, fatores externos ao comportamento estudantil podem influenciar múltiplas métricas simultaneamente, como instabilidades na conexão, limitações de \textit{hardware} e configuração do navegador. Esses elementos podem introduzir vieses que não refletem diferenças reais de engajamento, reforçando a necessidade de interpretar os resultados com cautela e em conjunto com dados qualitativos e contexto pedagógico.
